\begin{description}
\item[(D02.1a)] Average acceleration (in m\,s$^{-2}$) is given by the
  following relation,
  \[
  a = (v-u)\times 10^{3}/(\tau\times 3600)
  \]
  \begin{center}
  \begin{tabular}{l c c c c}
  \toprule
  CME & $u$ & $v$ & $\tau$ & $a$\\
  Name & (km\,s$^{-1}$) & (km\,s$^{-1}$) & (h) & (m\,s$^{-2}$)\\
  \midrule
  CME-A & 804  & 470 & 74.5  & $-1.25$\\
  CME-B & 247  & 360 & 127.5 & 0.246\\
  CME-C & 523  & 396 & 103.5 & $-0.341$\\
  CME-D & 830  & 415 & 71.0  & $-1.62$\\
  CME-E & 665  & 400 & 104.5 & $-0.704$\\
  CME-F & 347  & 350 & 101.5 & 0.00821\\
  CME-G & 446  & 375 & 99.5  & $-0.198$\\
  CME-H & 155  & 360 & 97.0  & 0.587\\
  CME-I & 1016 & 515 & 67.0  & $-2.08$\\
  CME-J & 683  & 410 & 54.0  & $-1.40$\\
  \bottomrule
  \end{tabular}
  \end{center}
\item[(D02.1b)]
  \begin{align*}
  m &= (-3.05 \pm 0.09)~\mathrm{m\,s^{-2}}\\
  \alpha &= (1.08 \pm 0.06)~\mathrm{m\,s^{-2}}
  \end{align*}
  Thus, the expression obtained by fitting a straight line to $a$ vs $u/u_0$
  is
  \[
  a_{\mathrm{model}} = -3.05\left(\frac{u}{u_0}\right) + 1.08
  \]
  % FIGURE NEEDED: straight-line fit graph "D02.1b" of a vs u/u0 (drawn by
  % the student; no model graph is printed in the official solution).
\item[(D02.1c)]
  \begin{center}
  \begin{tabular}{l c}
  \toprule
  CME & $a_{\mathrm{model}}$\\
  Name & (m\,s$^{-2}$)\\
  \midrule
  CME-A & $-1.37$\\
  CME-B & 0.327\\
  CME-C & $-0.515$\\
  CME-D & $-1.45$\\
  CME-E & $-0.948$\\
  CME-F & 0.0217\\
  CME-G & $-0.280$\\
  CME-H & 0.607\\
  CME-I & $-2.02$\\
  CME-J & $-1.00$\\
  \bottomrule
  \end{tabular}
  \end{center}
  The root mean square (rms) deviation of acceleration is given by
  \[
  \delta a_{\mathrm{rms}}
    = \sqrt{\frac{\sum_{i=1}^{10}(a_{\mathrm{model},i}-a_i)^2}{10}}
  \]
  Using this expression we get, $\delta a_{\mathrm{rms}} = 0.176$\,m\,s$^{-2}$.
\item[(D02.2a)] The time of arrival of a CME at the Earth, $\tau_m$, can be
  calculated from the following relation.
  \[
  s = u\tau_m + \frac{1}{2}a_{\mathrm{model}}\tau_m^2
  \]
  However, this would yield two solutions for $\tau_m$, only one of which will
  be for a positive final velocity (acceptable), and the other for a negative
  final velocity (not acceptable). Since acceleration is constant, the final
  velocity $v_m$ at Earth at a distance $s = 214\,R_\odot$ is given by
  \[
  v_m = \left(u^2 + 2a_{\mathrm{model}}s\right)^{1/2}
      = \left[u^2 + 2\left(m\frac{u}{u_0}+\alpha\right)s\right]^{1/2}
  \]
  Using this velocity, we get the arrival time $\tau$ as
  \[
  \tau_m = \frac{v_m - u}{a_{\mathrm{model}}}
         = \frac{v_m - u}{m\frac{u}{u_0}+\alpha}
  \]
  Using the given value of $u$, and the obtained values of $m$ and $\alpha$,
  we get

  For CME-1: $v_{1,m} = 680.728$\,km\,s$^{-1}$ which gives
  $\tau_{1,m} = 48.0$\,h.

  For CME-2: $v_{1,m} = 384.941$\,km\,s$^{-1}$ which gives
  $\tau_{2,m} = 126$\,h.
  % sic: the official solution writes v_{1,m} for CME-2 as well (should be
  % v_{2,m})
\item[(D02.2b)] The arrival time for CME-1 is $\tau_{1,m} = 48.0$\,h,
  therefore the percentage error in the arrival time of CME-1 is,
  \[
  \Delta\tau_{1,m} = \frac{48.0-46.0}{46.0}\times 100\% = 4.35\%
  \]
  Similarly, the arrival time for CME-2 is $\tau_{2,m} = 126$\,h, hence the
  percentage error is arrival time for CME-2 is,
  % sic: "error is arrival time" as in the original
  \[
  \Delta\tau_{2,m} = \frac{126-74.5}{74.5}\times 100\% = 69.1\%
  \]
  Hence CME-1: VALID; CME-2: NOT VALID.
\item[(D02.3a)] Inserting $R_0$, $t_0$, $V_0$, and $V_s$ in the given
  expression
  \[
  R_D(t) = \frac{S}{\gamma}\ln\left[1 + S\gamma(V_0-V_s)(t-t_0)\right]
         + V_s(t-t_0) + R_0,
  \]
  we get

  Case C1:
  \begin{align*}
  R_0 &= \frac{7.99+6.36}{2} = 7.18~R_\odot\\
  t_0 &= \frac{0.20+0.48}{2} = 0.340~\mathrm{h}\\
  V_0 &= \frac{(7.99-6.36)\times 695700}{(0.48-0.20)\times 3600}
       = 1.12\times 10^{3}~\mathrm{km\,s^{-1}}
  \end{align*}
  we get, $R_D = 210.61\,R_\odot$ at $t = 58.05$\,h. Thus,
  $\delta R_D = 214 - 210.61 = 3.39\,R_\odot$.

  Case C2:
  \begin{align*}
  R_0 &= \frac{11.99+13.51}{2} = 12.8~R_\odot\\
  t_0 &= \frac{1.22+1.49}{2} = 1.36~\mathrm{h}\\
  V_0 &= \frac{(13.51-11.99)\times 695700}{(1.49-1.22)\times 3600}
       = 1.09\times 10^{3}~\mathrm{km\,s^{-1}}
  \end{align*}
  Using the expression above, $R_D = 210.86\,R_\odot$ at $t=58.05$. Thus,
  $\delta R_D = 214 - 210.86 = 3.14\,R_\odot$.
\item[(D02.3b)] Case C3:
  \begin{align*}
  R_0 &= \frac{6+4}{2} = 5~R_\odot\\
  t_0 &= \frac{3+1}{2} = 2~\mathrm{h}\\
  V_0 &= \frac{(6-4)\times 695700}{(3-1)\times 3600} = 193~\mathrm{km\,s^{-1}}
  \end{align*}
  Case C4:
  \begin{align*}
  R_0 &= \frac{11+9}{2} = 10~R_\odot\\
  t_0 &= \frac{5+4}{2} = 4.5~\mathrm{h}\\
  V_0 &= \frac{(11-9)\times 695700}{(5-4)\times 3600} = 386~\mathrm{km\,s^{-1}}
  \end{align*}
  \begin{center}
  \begin{tabular}{c c c c}
  \toprule
  Data Point & $t$ (in h) & C3 & C4\\
  & & $R_D(t)$ (in $R_\odot$) & $R_D(t)$ (in $R_\odot$)\\
  \midrule
  P5 & 21.0 & 25.1 & 42.8\\
  P6 & 50.0 & 59.1 & 99.9\\
  P7 & 85.0 & 104  & 168\\
  P8 & 111  & 139  & 218\\
  \bottomrule
  \end{tabular}
  \end{center}
\item[(D02.3c)] The plot below shows the data points of CME-4 and smooth
  variation of $R_D$ extrapolated up to $220\,R_\odot$.

  % FIGURE NEEDED: graph "D02.3c" -- R_D(t) (in R_sun) vs t (in hours) for
  % cases C3 and C4 of CME-4, x axis 0--180 h, with smooth curves through
  % P5--P8; 2025_data_analysis_solutions.pdf, p.~21.
\item[(D02.3d)] The arrival time of CME at Earth ($214\,R_\odot$) can be
  obtained by estimating the point of intersection between
  $R_D = 214\,R_\odot = $ constant and two smooth curves.

  The points of intersection are 167\,h and 109\,h for the cases C3 and C4,
  respectively. Thus, absolute time differences are:
  \[
  \text{C3: } |\delta\tau| = 56~\mathrm{h}
  \qquad\text{and}\qquad
  \text{C4: } |\delta\tau| = 2~\mathrm{h}.
  \]
\item[(D02.3e)] The statement is TRUE.

  In the scenario of CME-3, the ``drag-only'' model predicts the arrival
  distance accurately when $R_D$ is computed using the initial velocity,
  $V_0$ and initial distance, $R_0$ at 0.20\,h and beyond. While for CME-4,
  the ``drag-only'' model start accurately predicting the arrival times of
  CME-4 when $R_D$ is calculated using $V_0$ and $R_0$ at 4\,h after the
  launch of the CME.
  % sic: "model start" as in the original
\item[(D02.4)]
  \[
  V_D(t) = \frac{V_0 - V_s}{1 + S\gamma(V_0-V_s)(t-t_0)} + V_s
  \qquad\Longrightarrow\qquad
  v = \frac{u - V_s}{1 + S\gamma(u-V_s)\tau} + V_s.
  \]
  Taking, $S = 1$ if $v < u$, else $S = -1$. The above expression is
  quadratic equation, thus, we get two values of $V_s$. If $v < u$, choose
  the solution with $V_s < u$. If $v > u$, choose the solution with
  $V_s > u$.

  Calculating $V_s$ using the above conditions for all ten CMEs.
  \begin{center}
  \begin{tabular}{l c c c c}
  \toprule
  CME & $u$ & $v$ & $\tau$ & $V_s$\\
  Name & (km\,s$^{-1}$) & (km\,s$^{-1}$) & (h) & (km\,s$^{-1}$)\\
  \midrule
  CME-A & 804  & 470 & 74.5  & 337\\
  CME-B & 247  & 360 & 127.5 & 428\\
  CME-C & 523  & 396 & 103.5 & 314\\
  CME-D & 830  & 415 & 71.0  & 270\\
  CME-E & 665  & 400 & 104.5 & 303\\
  CME-F & 347  & 350 & 101.5 & 369\\
  CME-G & 446  & 375 & 99.5  & 305\\
  CME-H & 155  & 360 & 97.0  & 457\\
  CME-I & 1016 & 515 & 67.0  & 357\\
  CME-J & 683  & 410 & 54.0  & 248\\
  \bottomrule
  \end{tabular}
  \end{center}
  Taking the average of the solar wind speed, $V_s$, for each case, we get
  $V_{s,\mathrm{avg}} = 339$\,km\,s$^{-1}$.
\end{description}
